\documentclass[a4paper,12pt]{article}

\usepackage[italian]{babel}
\usepackage[T1]{fontenc}
\usepackage[utf8]{inputenc}
\usepackage{lmodern}
\usepackage{geometry}
\usepackage{amsmath}
\usepackage{booktabs}
\usepackage{graphicx}
\usepackage{float}
\usepackage{hyperref}
\usepackage{siunitx}
\usepackage{enumitem}

\geometry{margin=2.5cm}
\hypersetup{colorlinks=true,linkcolor=blue,urlcolor=blue,citecolor=blue}

\title{Analisi morfologica e classificazione non-CNN di lettere manoscritte}
\author{Studente: \rule{6cm}{0.4pt}}
\date{\today}

\begin{document}

\maketitle

\begin{abstract}
Questo lavoro propone una pipeline completa e interpretabile per il riconoscimento di lettere manoscritte basata su Morfologia Matematica.
L'idea centrale è separare in modo esplicito le fasi di: (i) pulizia dell'immagine, (ii) estrazione di struttura geometrica/topologica, (iii) classificazione supervisionata.
Le immagini sono binarizzate (Otsu), scheletrizzate e potate (pruning) per ridurre il rumore mantenendo l'informazione di forma.
Da queste rappresentazioni vengono estratte feature robuste (topologiche, geometriche, spaziali, statistiche), poi usate da modelli non-CNN (SVM RBF, ExtraTrees, XGBoost).
La valutazione è svolta con split per soggetto e con cross-validation subject-wise, per misurare la reale capacità di generalizzare a calligrafie mai viste.
Il miglior risultato holdout è ottenuto da XGBoost (Accuracy 0.7369, Macro-F1 0.7338), confermando che il feature engineering morfologico è competitivo e ben motivato.
\end{abstract}

\section{Introduzione}
Il riconoscimento di scrittura manoscritta è un problema complesso per tre motivi principali:
\begin{itemize}[leftmargin=1.5em]
  \item \textbf{variabilità inter-soggetto}: la stessa lettera può avere forme molto diverse;
  \item \textbf{rumore di acquisizione}: sfondo, spessori irregolari, piccole sbavature;
  \item \textbf{ambiguità intrinseca}: alcune coppie di lettere sono naturalmente simili.
\end{itemize}

Invece di una soluzione end-to-end con rete neurale profonda, questo progetto adotta un approccio \textbf{esplicativo}: ogni passaggio è leggibile e ogni feature ha un significato geometrico o topologico chiaro. 

\paragraph{Obiettivi.}
\begin{itemize}[leftmargin=1.5em]
  \item Ridurre il rumore preservando la struttura del tratto.
  \item Estrarre descrittori topologici e geometrici robusti.
  \item Validare la generalizzazione tra soggetti diversi con protocollo subject-wise.
  \item Mantenere una pipeline riproducibile e interpretabile in tutte le fasi.
\end{itemize}

\section{Dataset e organizzazione}
Il dataset è organizzato per soggetto e classe, nella forma:
\begin{center}
\texttt{data/lettere/<soggetto>/<lettera>/<immagine>}
\end{center}

Questa struttura permette di tenere separata l'identità del soggetto dalla classe della lettera, condizione fondamentale per evitare data leakage quando si fa split train/test.

Statistiche principali:
\begin{itemize}[leftmargin=1.5em]
  \item \textbf{14521} campioni totali,
  \item \textbf{25} classi di lettere,
  \item \textbf{20} soggetti.
\end{itemize}

In pratica, ogni campione attraversa la stessa pipeline morfologica e viene trasformato in un vettore numerico uniforme, poi aggregabile per analisi globale o per soggetto.

\subsection{Preparazione del dataset: estrazione caratteri da CVL}
La prima fase del progetto non è la classificazione, ma la costruzione del dataset a partire dalle pagine manoscritte.
Per questo è stato aggiunto uno script dedicato, \texttt{extract\_characters\_cvl.py}, che trasforma una pagina completa in singoli crop di carattere pronti per l'analisi successiva.

In altre parole, questa fase fa da ponte tra i manoscritti grezzi e la pipeline morfologica: senza segmentazione dei caratteri, non avremmo campioni omogenei su cui calcolare le feature e addestrare i modelli.

Lo script opera su immagine singola o in batch e realizza:
\begin{enumerate}[leftmargin=1.5em]
  \item conversione in grayscale,
  \item binarizzazione adattiva con soglia gaussiana e inversione del foreground,
  \item denoising per attenuare piccoli artefatti e rumore impulsivo,
  \item erosione e dilatazione per ripulire il tratto e stabilizzare i contorni,
  \item eventuale opening/closing per rimuovere micro-componenti isolate o chiudere piccole interruzioni,
  \item estrazione delle componenti connesse e delle bounding box candidate,
  \item filtraggio geometrico su area, larghezza, altezza e aspect ratio,
  \item filtro outlier basato su mediane (area/altezza) per scartare box troppo grandi rispetto al comportamento tipico della pagina,
  \item ordinamento in ordine di lettura (riga per riga, da sinistra a destra),
  \item salvataggio dei crop carattere per carattere,
  \item esportazione del file \texttt{extraction\_manifest.csv} con coordinate (x, y, w, h).
\end{enumerate}

Dal punto di vista operativo, quindi, l'immagine viene prima normalizzata e ripulita, poi segmentata in regioni candidate, e infine rifinita con controlli geometrici per ottenere crop di carattere più consistenti e utili alle fasi successive.

Questa fase è coerente con il progetto perché rende esplicito il passaggio da pagina manoscritta a campioni elementari (caratteri) su cui poi si applica la pipeline morfologica e la classificazione.

\paragraph{Configurazione consigliata (quality profile).}
Per ottenere una segmentazione più pulita sui manoscritti CVL, la configurazione consigliata è:
\begin{itemize}[leftmargin=1.5em]
  \item \texttt{--min-area 30}
  \item \texttt{--min-width 5}
  \item \texttt{--min-height 8}
  \item \texttt{--max-aspect-ratio 10}
  \item \texttt{--opening-kernel 2}
\end{itemize}
Questa scelta riduce in modo sensibile micro-componenti spurie e rettangoli anomali, mantenendo i caratteri utili per le fasi successive di feature extraction e classificazione.

Dal punto di vista del flusso complessivo, questa è la parte che precede tutto il resto: prima estraggo i caratteri, poi applico la morfologia, poi costruisco le feature e infine alleno i modelli.

\subsection{Demo visiva della pipeline}
Per illustrare passo-passo l'intero flusso è stato aggiunto un modulo di visualizzazione: \texttt{pipeline\_visual\_classifier.py}. 
Insieme a \texttt{train\_xgboost\_bundle.py} (che allena un modello XGBoost e salva un bundle contenente modello, imputer e label-encoder), il modulo di visualizzazione esegue: estrazione dei box, salvataggio di un'immagine di preprocessing della pagina, generazione di pagine-grid con i crop classificati e figure di dettaglio per singoli caratteri.

Questa è la parte che puoi mostrare subito quando ti chiedono di vedere i risultati in modo visivo: la sequenza va dalla pagina intera ai crop, fino alla classificazione del singolo carattere.

Per eseguire la demo (dopo aver generato il bundle con \texttt{train\_xgboost\_bundle.py}):
\begin{verbatim}
python pipeline_visual_classifier.py --input data/manoscritti/0002-1.tif --bundle-path output/classification_advanced/organized/06_model_artifacts/xgboost_bundle.joblib --output output/pipeline_demo
\end{verbatim}

Il modulo genera una cartella di report, ad esempio `output/pipeline_demo/0002-1/`, contenente immagini e file utili alla presentazione.

Esempi di immagini create dalla demo, in ordine di lettura della pipeline:
\begin{figure}[H]
  \centering
  \includegraphics[width=0.85\textwidth]{output/pipeline_demo/0002-1/01_page_processing.png}
  \caption{Elaborazione della pagina - documento originale, grayscale, binarizzazione e bounding box rilevati.\newline
  \textit{Cosa mostra:} i principali step di preprocessing su una pagina di manoscritto; \textit{Cosa ci dice:} come vengono individuati i candidati crop e quanto pulita è la binarizzazione prima di estrarre i caratteri.}
\end{figure}

\begin{figure}[H]
  \centering
  \includegraphics[width=0.85\textwidth]{output/pipeline_demo/0002-1/02_extracted_characters.png}
  \caption{Panoramica dei caratteri estratti e classificati.\newline
  \textit{Cosa mostra:} i crop ottenuti dalla segmentazione, già accompagnati dalla lettera predetta e dalla confidenza del modello; \textit{Cosa ci dice:} permette di capire al volo se l'estrazione è pulita e se il classificatore sta leggendo correttamente i caratteri.}
\end{figure}

\begin{figure}[H]
  \centering
  \includegraphics[width=0.85\textwidth]{output/pipeline_demo/0002-1/03_selected_character_detail.png}
  \caption{Dettaglio di un carattere selezionato: crop, binarizzazione e scheletro pulito con feature principali.\newline
  \textit{Cosa mostra:} visualizzazione affiancata del crop, la sua binarizzazione e lo scheletro dopo pruning, con le feature numeriche più rilevanti; \textit{Cosa ci dice:} permette di spiegare la predizione caso per caso e verificare se errori sono causati dall'estrazione o dalla rappresentazione feature.}
\end{figure}

\begin{figure}[H]
  \centering
  \includegraphics[width=0.85\textwidth]{output/pipeline_demo/0002-1/04_character_details_batch/04_character_detail_0001_char_001_a.png}
  \caption{Esempio di figure dettagliate generate per tutti i caratteri (in cartelle per lettera).\newline
  \textit{Cosa mostra:} singoli esempi salvati per lettera con valore predetto e principali feature; \textit{Cosa ci dice:} utile per analisi errori sistematici e per selezionare esempi per la presentazione o per l'annotazione umana.}
\end{figure}

Queste immagini sono pensate per essere incluse nella presentazione: mostrano l'intera catena operativa e permettono di motivare tutte le scelte di preprocessing, estrazione e classificazione.

\section{Metodologia}
\subsection{Preprocessing e morfologia}
La pipeline operativa include quattro passaggi principali:
\begin{enumerate}[leftmargin=1.5em]
  \item \textbf{Conversione e binarizzazione (Otsu).} L'immagine in scala di grigi viene separata in foreground/sfondo con una soglia automatica. Questo rende robusta la segmentazione del tratto senza fissare soglie manuali.
  \item \textbf{Skeletonization.} Il tratto binario viene ridotto a uno scheletro a 1 pixel di spessore, preservando connettività e struttura globale.
  \item \textbf{Pruning iterativo.} Si rimuovono iterativamente endpoint troppo corti (rami spurii), in modo da attenuare rumore locale e micro-ramificazioni non informative.
  \item \textbf{Estrazione feature.} Le misure sono calcolate sia sull'immagine binaria sia sullo scheletro ripulito, così da catturare informazioni complementari.
\end{enumerate}

\subsection{Perché la morfologia è centrale}
Le operazioni morfologiche non sono un semplice preprocessing ``estetico'': costituiscono la base semantica del modello.
In particolare:
\begin{itemize}[leftmargin=1.5em]
  \item la binarizzazione definisce il dominio su cui misurare area, densità e cavità;
  \item lo scheletro rende esplicita la topologia del carattere (terminazioni, incroci, lunghezza dei rami);
  \item il pruning aumenta il rapporto segnale/rumore delle feature strutturali.
\end{itemize}

Senza questi passaggi, molte feature perderebbero stabilità e interpretabilità.

\subsection{Feature estratte}
Le feature sono organizzate in gruppi complementari. Di seguito sono descritte con il loro significato operativo.

\paragraph{1) Feature base/topologiche}
\begin{itemize}[leftmargin=1.5em]
  \item \textbf{endpoints / punte}: numero di terminazioni dello scheletro; utile per distinguere lettere aperte/chiuse o con code.
  \item \textbf{junctions / incroci}: numero di nodi con più rami; discriminante per lettere con intersezioni interne.
  \item \textbf{holes / buchi}: numero di cavità chiuse nel tratto (es. ``o'', ``p'', ``b'').
  \item \textbf{componenti}: numero di componenti connesse, utile in presenza di tratti separati.
  \item \textbf{euler\_number}: sintetizza topologia globale (componenti meno buchi).
\end{itemize}

\paragraph{2) Feature geometriche di base}
\begin{itemize}[leftmargin=1.5em]
  \item \textbf{aspect\_ratio}: rapporto larghezza/altezza della bounding box della componente principale.
  \item \textbf{densita\_pixel}: frazione di pixel attivi sull'area totale.
  \item \textbf{avg\_stroke\_width}: stima dello spessore medio del tratto (distance transform + scheletro).
  \item \textbf{skeleton\_length}: lunghezza complessiva del tratto scheletrizzato.
  \item \textbf{skeleton\_density}: densità dello scheletro rispetto all'ink totale.
  \item \textbf{endpoint\_norm, junction\_norm}: endpoint/incroci normalizzati per lunghezza scheletro.
  \item \textbf{pruning\_removed\_ratio}: quota di scheletro rimossa dal pruning (indicatore di rumore locale).
\end{itemize}

\paragraph{3) Feature di forma globale e invarianti}
\begin{itemize}[leftmargin=1.5em]
  \item \textbf{solidity}: rapporto area/convex hull area, misura ``pienezza'' della forma.
  \item \textbf{extent}: rapporto area/area bounding box.
  \item \textbf{eccentricity}: allungamento della forma principale.
  \item \textbf{orientation}: orientazione prevalente dell'asse maggiore.
  \item \textbf{major\_axis\_length, minor\_axis\_length, axis\_ratio}: parametri ellittici della componente principale.
  \item \textbf{hu\_1 ... hu\_7}: momenti invarianti di Hu (log-compressi), robusti a traslazione/scala/rotazione.
\end{itemize}

\paragraph{4) Feature di struttura del grafo scheletrico}
\begin{itemize}[leftmargin=1.5em]
  \item \textbf{n\_branches}: numero di rami ottenuti rimuovendo i nodi di giunzione.
  \item \textbf{branch\_length\_mean, branch\_length\_max, branch\_length\_std}: statistiche di lunghezza dei rami.
\end{itemize}

\paragraph{5) Feature spaziali}
\begin{itemize}[leftmargin=1.5em]
  \item \textbf{zoning 3x3} (\texttt{zone\_r*\_c*\_density}): densità locale in ciascuna cella della griglia.
  \item \textbf{projection profiles}: statistiche dei profili orizzontali e verticali.
  \item \textbf{spatial balance}: bilanciamento top/bottom e left/right (es. \texttt{top\_bottom\_ratio}, \texttt{left\_right\_ratio}).
\end{itemize}

\paragraph{6) Feature strutturali avanzate}
\begin{itemize}[leftmargin=1.5em]
  \item \textbf{transition features} (riga/colonna): numero e statistiche di transizioni 0↔1, indicano complessità del contorno.
  \item \textbf{diagonal profile}: densità su diagonale principale e secondaria, con differenza tra le due.
  \item \textbf{endpoint position}: posizione media/dispersione normalizzata delle terminazioni.
\end{itemize}

\paragraph{7) Feature confusion-driven}
Queste feature sono pensate per migliorare coppie notoriamente ambigue:
\begin{itemize}[leftmargin=1.5em]
  \item \textbf{scanline structure}: pattern locali su righe campione;
  \item \textbf{border contact}: come il tratto tocca i bordi dell'immagine normalizzata;
  \item \textbf{hole position}: posizione dei buchi (non solo quantità);
  \item \textbf{contour depth}: profondità media del tratto verso l'interno;
  \item \textbf{endpoint distribution}: distribuzione spaziale delle terminazioni;
  \item \textbf{midline runs}: lunghezza dei segmenti continui su linee mediane.
\end{itemize}

\paragraph{8) Statistiche aggiuntive sulle proiezioni}
Per \texttt{row\_proj} e \texttt{col\_proj} vengono aggiunti:
\begin{itemize}[leftmargin=1.5em]
  \item quantili ($q_{25}, q_{50}, q_{75}$),
  \item asimmetria (skewness),
  \item curtosi (kurtosis).
\end{itemize}
Queste misure migliorano la descrizione della distribuzione del tratto oltre la sola media/varianza.

\subsection{Spiegazione esplicita del \textit{perché} di ogni feature}
In questa sezione ogni famiglia è collegata direttamente al tipo di errore che aiuta a ridurre.

\paragraph{Topologia (\texttt{punte}, \texttt{incroci}, \texttt{buchi}, \texttt{componenti}, \texttt{euler\_number}).}
\begin{itemize}[leftmargin=1.5em]
  \item \texttt{punte}: distingue lettere con code/terminazioni nette (es. tratti finali aperti) da forme più chiuse.
  \item \texttt{incroci}: separa lettere con struttura ramificata da lettere ``a tratto unico''.
  \item \texttt{buchi}: discrimina lettere con anello interno da lettere senza cavità.
  \item \texttt{componenti}: identifica casi con tratti disconnessi (utile in presenza di scrittura spezzata o rumore).
  \item \texttt{euler\_number}: sintetizza in un solo valore la relazione componenti/buchi, rendendo più robusta la distinzione topologica.
\end{itemize}

\paragraph{Geometria globale (\texttt{aspect\_ratio}, \texttt{densita\_pixel}, \texttt{avg\_stroke\_width}, \texttt{skeleton\_length}, \texttt{skeleton\_density}, \texttt{endpoint\_norm}, \texttt{junction\_norm}, \texttt{pruning\_removed\_ratio}).}
\begin{itemize}[leftmargin=1.5em]
  \item \texttt{aspect\_ratio}: separa lettere tendenzialmente ``alte'' da lettere ``larghe''.
  \item \texttt{densita\_pixel}: misura quanta ``inchiostrazione'' c'è; utile per differenze di riempimento del gesto.
  \item \texttt{avg\_stroke\_width}: cattura lo stile di pressione/spessore del tratto.
  \item \texttt{skeleton\_length}: proxy della complessità/estensione del percorso scritto.
  \item \texttt{skeleton\_density}: confronta lunghezza scheletro e area inchiostrata, distinguendo tratti sottili da forme compatte.
  \item \texttt{endpoint\_norm}, \texttt{junction\_norm}: rendono confrontabili lettere di dimensione diversa grazie alla normalizzazione.
  \item \texttt{pruning\_removed\_ratio}: indica quanto rumore ramificato era presente; è informativa per stili tremolanti o sbavati.
\end{itemize}

\paragraph{Forma principale (\texttt{solidity}, \texttt{extent}, \texttt{eccentricity}, \texttt{orientation}, \texttt{major\_axis\_length}, \texttt{minor\_axis\_length}, \texttt{axis\_ratio}).}
\begin{itemize}[leftmargin=1.5em]
  \item \texttt{solidity}: distingue forme compatte da forme con rientranze.
  \item \texttt{extent}: misura quanto la lettera occupa la propria bounding box.
  \item \texttt{eccentricity} e \texttt{axis\_ratio}: misurano l'allungamento e aiutano su coppie con simile topologia ma proporzioni diverse.
  \item \texttt{orientation}: cattura inclinazione prevalente (utile su calligrafie inclinate).
  \item \texttt{major/minor\_axis\_length}: aggiungono scala e forma ellittica della componente dominante.
\end{itemize}

\paragraph{Invarianti di momento (\texttt{hu\_1} ... \texttt{hu\_7}).}
Le Hu moments sono utili perché conservano informazione di forma anche in presenza di piccole variazioni geometriche; riducono la sensibilità a traslazione/scala/rotazione e quindi migliorano la robustezza inter-soggetto.

\paragraph{Struttura dei rami (\texttt{n\_branches}, \texttt{branch\_length\_mean}, \texttt{branch\_length\_max}, \texttt{branch\_length\_std}).}
Queste feature descrivono la ``grammatica'' interna dello scheletro: quante diramazioni ci sono e quanto sono lunghe/variabili. Sono particolarmente utili quando due lettere hanno area simile ma struttura interna differente.

\paragraph{Distribuzione spaziale locale (\texttt{zone\_r*\_c*\_density}, \texttt{row\_proj\_*}, \texttt{col\_proj\_*}, \texttt{top\_bottom\_ratio}, \texttt{left\_right\_ratio}).}
\begin{itemize}[leftmargin=1.5em]
  \item lo \textbf{zoning 3x3} localizza dove si concentra il tratto;
  \item i \textbf{projection profiles} descrivono il profilo verticale/orizzontale della forma;
  \item il \textbf{spatial balance} separa lettere ``sbilanciate in alto/basso'' o ``a sinistra/destra''.
\end{itemize}
Insieme riducono confusioni tra lettere con stessa area totale ma diversa disposizione del tratto.

\paragraph{Complessità del contorno (\texttt{row\_trans\_*}, \texttt{col\_trans\_*}, \texttt{main\_diag\_density}, \texttt{anti\_diag\_density}, \texttt{diag\_density\_diff}).}
Le transizioni 0↔1 quantificano quante volte il segnale entra/esce dal tratto lungo righe/colonne: più transizioni implicano contorni più complessi. Le diagonali aggiungono sensibilità alla direzione dominante del gesto, utile per lettere inclinate o oblique.

\paragraph{Posizionamento delle terminazioni (\texttt{endpoint\_x\_mean\_norm}, \texttt{endpoint\_y\_mean\_norm}, \texttt{endpoint\_x\_std\_norm}, \texttt{endpoint\_y\_std\_norm}, \texttt{endpoint\_distribution\_*}).}
Queste misure spiegano \textit{dove} finiscono i tratti: terminazioni concentrate in alto, in basso o distribuite su tutta la lettera. Sono utili per separare lettere con stesso numero di endpoint ma disposizione differente.

\paragraph{Feature confusion-driven (\texttt{scanline\_*}, \texttt{border\_contact\_*}, \texttt{hole\_position\_*}, \texttt{contour\_depth\_*}, \texttt{midline\_run\_*}).}
Sono state introdotte per migliorare coppie ambigue osservate in confusion matrix.
\begin{itemize}[leftmargin=1.5em]
  \item \texttt{scanline\_*}: confrontano pattern locali su linee specifiche.
  \item \texttt{border\_contact\_*}: rilevano se il tratto tocca certi bordi (informazione discriminativa di forma aperta/chiusa).
  \item \texttt{hole\_position\_*}: aggiungono la posizione delle cavità, non solo il conteggio.
  \item \texttt{contour\_depth\_*}: misurano ``profondità'' e penetrazione del tratto rispetto al bordo.
  \item \texttt{midline\_run\_*}: quantificano segmenti continui sulle linee mediane, utili per distinguere andamenti centrali del tratto.
\end{itemize}

\paragraph{Statistiche di forma dei profili (\texttt{row/col\_proj\_q25/q50/q75}, \texttt{row/col\_proj\_skew}, \texttt{row/col\_proj\_kurtosis}).}
Queste statistiche non descrivono solo ``quanto'' tratto è presente, ma \textit{come} è distribuito: asimmetria e curtosi aiutano a distinguere profili con stessa media ma forma diversa.

\section{Setup sperimentale}
\subsection{Split holdout per soggetto}
Per testare la generalizzazione reale, la divisione è fatta sui soggetti e non sui singoli campioni.
Questo significa che il modello in test vede stili grafici completamente nuovi.

\begin{itemize}[leftmargin=1.5em]
  \item train: 15 soggetti (11610 campioni)
  \item test: 5 soggetti (2911 campioni)
\end{itemize}

\subsection{Cross-validation subject-wise}
Per robustezza statistica è usata anche una CV a 5 fold subject-wise.
Ogni fold mantiene la separazione per soggetto, così la varianza tra fold riflette differenze di stile realistiche.

\subsection{Modelli e metriche}
Sono confrontati tre modelli con ipotesi diverse:
\begin{itemize}[leftmargin=1.5em]
  \item \textbf{SVM RBF}: adatto a separazioni non lineari in spazi ad alta dimensione;
  \item \textbf{ExtraTrees}: ensemble ad albero molto robusto a feature eterogenee;
  \item \textbf{XGBoost}: boosting gradient-based, spesso molto efficace su feature engineering tabellare.
\end{itemize}

Metriche principali:
\begin{itemize}[leftmargin=1.5em]
  \item \textbf{Accuracy}: quota complessiva di predizioni corrette;
  \item \textbf{Macro-F1}: media F1 per classe con peso uniforme (fondamentale in presenza di difficoltà diverse tra lettere).
\end{itemize}

Completano l'analisi: report per classe, confusion matrix, F1 per lettera.

\section{Risultati}
\subsection{Risultati holdout aggiornati}
\begin{table}[H]
\centering
\caption{Confronto modelli su holdout subject-wise}
\begin{tabular}{lcc}
\toprule
Modello & Accuracy & Macro-F1 \\
\midrule
XGBoost & 0.7369 & 0.7338 \\
SVM RBF & 0.7248 & 0.7239 \\
ExtraTrees & 0.7238 & 0.7205 \\
\bottomrule
\end{tabular}
\end{table}

Il miglior modello in holdout risulta \textbf{XGBoost}.
La differenza, pur non enorme, è coerente su entrambe le metriche principali e indica una migliore capacità di sfruttare la varietà delle feature estratte.

\subsection{Risultati CV subject-wise (riferimento)}
\begin{table}[H]
\centering
\caption{CV 5-fold subject-wise (media $\pm$ deviazione standard)}
\begin{tabular}{lcc}
\toprule
Modello & Accuracy & Macro-F1 \\
\midrule
XGBoost & $0.7676 \pm 0.0246$ & $0.7236 \pm 0.0269$ \\
SVM RBF & $0.7372 \pm 0.0427$ & $0.7039 \pm 0.0351$ \\
ExtraTrees & $0.7472 \pm 0.0265$ & $0.6898 \pm 0.0306$ \\
\bottomrule
\end{tabular}
\end{table}

In CV si osserva una buona stabilità del modello migliore: la deviazione standard resta contenuta e conferma che il comportamento non dipende da un singolo split favorevole.


\subsection{Ablation study dei gruppi di feature}
Per quantificare il contributo dei diversi blocchi di feature è stato eseguito un ablation study con XGBoost: si rimuove un gruppo alla volta e si misura la variazione di performance rispetto alla configurazione \texttt{all\_features}.

Nel setup di riferimento dell'ablation (stesso split subject-wise in cache) la baseline è:
\begin{itemize}[leftmargin=1.5em]
  \item Accuracy = 0.7269
  \item Macro-F1 = 0.7250
\end{itemize}

\begin{figure}[H]
  \centering
  \includegraphics[width=0.9\textwidth]{ablation_delta_macro_f1.png}
  \caption{Ablation XGBoost: variazione di Macro-F1 quando si rimuove ciascun gruppo di feature.\newline
  \textit{Cosa mostra:} l'effetto marginale di ogni gruppo sulla qualità del modello. \textit{Cosa ci dice:} le barre negative identificano gruppi informativi, le positive gruppi potenzialmente ridondanti nel setup corrente.}
  \label{fig:ablation_delta}
\end{figure}

L'immagine in Figura~\ref{fig:ablation_delta} si interpreta così:
\begin{itemize}[leftmargin=1.5em]
  \item barre a sinistra dello zero (\(\Delta\) Macro-F1 negativo) indicano gruppi \textbf{importanti}: rimuoverli peggiora il modello;
  \item barre a destra dello zero (\(\Delta\) Macro-F1 positivo) indicano gruppi potenzialmente \textbf{ridondanti} o rumorosi in questo setup.
\end{itemize}

I risultati più rilevanti sono riportati in Tabella~\ref{tab:ablation_key_groups}.

\begin{table}[H]
\centering
\caption{Ablation XGBoost: gruppi principali e impatto su Macro-F1.}
\label{tab:ablation_key_groups}
\begin{tabular}{lc}
\toprule
Gruppo rimosso & $\Delta$ Macro-F1 vs all\_features \\
\midrule
\texttt{projection\_profiles} & -0.0166 \\
\texttt{transition\_profiles} & -0.0063 \\
\texttt{engineered\_interactions} & -0.0039 \\
\texttt{global\_shape} & -0.0028 \\
\texttt{stroke\_geometry} & +0.0119 \\
\texttt{hu\_moments} & +0.0101 \\
\bottomrule
\end{tabular}
\end{table}

Il motivo della perdita è che questi gruppi non descrivono dettagli secondari, ma aspetti strutturali molto informativi della lettera.
\texttt{projection\_profiles} cattura come l'inchiostro si distribuisce lungo righe e colonne, quindi separa bene lettere con massa concentrata in zone diverse; quando lo rimuovo, il modello perde informazione sulla geometria globale.
\texttt{transition\_profiles} misura la complessità locale del contorno e la frequenza di passaggi tra fondo e tratto, perciò aiuta a distinguere lettere con profili simili ma bordi differenti.
\texttt{engineered\_interactions} combina feature base già utili in forma non lineare, quindi la loro rimozione elimina relazioni tra endpoint, buchi e densità che XGBoost sfrutta molto bene.
\texttt{global\_shape} racchiude invece proprietà di forma complessiva come solidità, eccentricità e aspect ratio, cioè segnali importanti per distinguere lettere verticali, arrotondate o allungate.

Al contrario, il miglioramento leggero ottenuto togliendo \texttt{stroke\_geometry} e \texttt{hu\_moments} suggerisce che, nello split corrente, queste feature siano in parte ridondanti rispetto ad altri descrittori già presenti oppure più sensibili a variazioni di scrittura e normalizzazione.

In particolare, la rimozione di \texttt{projection\_profiles} produce la perdita più alta (\(-0.0166\) Macro-F1), quindi questo blocco risulta il più informativo nel setup corrente. Al contrario, la rimozione di \texttt{stroke\_geometry} e \texttt{hu\_moments} migliora leggermente la metrica in questo split: questo suggerisce possibile ridondanza e motiva una successiva feature selection mirata.


\section{Analisi qualitativa e diagnostica}
Per comprendere \textit{come} il modello decide e \textit{dove} sbaglia, l'analisi qualitativa include:
\begin{itemize}[leftmargin=1.5em]
  \item \textbf{heatmap correlazioni feature}: individua ridondanze e gruppi informativi;
  \item \textbf{proiezione PCA 2D}: visualizza la separazione globale tra classi;
  \item \textbf{confusion matrix}: mostra le coppie di lettere critiche;
  \item \textbf{F1 per lettera}: misura quali classi sono deboli/forti in modo puntuale.
\end{itemize}

\begin{figure}[H]
    \centering
    \includegraphics[width=0.8\textwidth]{feature_correlation_heatmap.png}
  \caption{Heatmap di correlazione (gruppi di feature).\newline
  \textit{Cosa mostra:} il livello di correlazione media tra blocchi di feature. \textit{Cosa ci dice:} blocchi molto correlati possono introdurre ridondanza, blocchi meno correlati aggiungono informazione complementare.}
\end{figure}

\begin{figure}[H]
    \centering
    \includegraphics[width=0.75\textwidth]{feature_space_pca_2d.png}
  \caption{PCA 2D del feature space.\newline
  \textit{Cosa mostra:} la proiezione bidimensionale delle lettere nello spazio delle feature. \textit{Cosa ci dice:} cluster separati indicano buona discriminabilità, sovrapposizioni anticipano possibili confusioni tra classi simili.}
\end{figure}

\begin{figure}[H]
    \centering
    \includegraphics[width=0.85\textwidth]{confusion_matrix_xgboost.png}
  \caption{Confusion matrix del modello XGBoost (holdout).\newline
  \textit{Cosa mostra:} conteggi assoluti di predizioni corrette/errate per ogni coppia vera-predetta. \textit{Cosa ci dice:} evidenzia dove il best model concentra gli errori e quali lettere risultano più stabili.}
\end{figure}

\begin{figure}[H]
  \centering
  \includegraphics[width=0.85\textwidth]{confusion_matrix_svm_rbf.png}
  \caption{Confusion matrix del modello SVM RBF (holdout).\newline
  \textit{Cosa mostra:} errori e successi assoluti della SVM sulle classi di test. \textit{Cosa ci dice:} permette il confronto diretto con XGBoost sulle stesse coppie di lettere difficili.}
\end{figure}

\begin{figure}[H]
  \centering
  \includegraphics[width=0.85\textwidth]{confusion_matrix_extra_trees.png}
  \caption{Confusion matrix del modello ExtraTrees (holdout).\newline
  \textit{Cosa mostra:} distribuzione degli errori in valore assoluto per il modello ad alberi randomizzati. \textit{Cosa ci dice:} chiarisce se il comportamento dell'ensemble differisce da SVM/XGBoost su lettere graficamente affini.}
\end{figure}

\begin{figure}[H]
  \centering
  \includegraphics[width=0.85\textwidth]{confusion_matrix_xgboost_normalized.png}
  \caption{Confusion matrix normalizzata del modello XGBoost.\newline
  \textit{Cosa mostra:} tassi di errore/correttezza per classe indipendenti dalla numerosità. \textit{Cosa ci dice:} consente di vedere quali lettere sono intrinsecamente più complesse da riconoscere anche con supporto diverso.}
\end{figure}

\begin{figure}[H]
  \centering
  \includegraphics[width=0.85\textwidth]{confusion_matrix_svm_rbf_normalized.png}
  \caption{Confusion matrix normalizzata del modello SVM RBF.\newline
  \textit{Cosa mostra:} probabilità di confusione classe-per-classe per la SVM. \textit{Cosa ci dice:} rende confrontabile la robustezza della SVM su classi rare e classi frequenti.}
\end{figure}

\begin{figure}[H]
  \centering
  \includegraphics[width=0.85\textwidth]{confusion_matrix_extra_trees_normalized.png}
  \caption{Confusion matrix normalizzata del modello ExtraTrees.\newline
  \textit{Cosa mostra:} profilo di errore relativo del modello ExtraTrees. \textit{Cosa ci dice:} aiuta a identificare quali classi perdono più precisione/recall rispetto agli altri modelli.}
\end{figure}

\begin{figure}[H]
  \centering
  \includegraphics[width=0.85\textwidth]{xgboost_feature_importance_top20.png}
  \caption{Feature importance (Top-20) del modello XGBoost.\newline
  \textit{Cosa mostra:} le feature con maggior contributo alle decisioni del boosting. \textit{Cosa ci dice:} evidenzia i segnali morfologici più discriminativi e guida eventuale feature selection.}
\end{figure}

\begin{figure}[H]
    \centering
    \includegraphics[width=0.85\textwidth]{f1_by_letter_xgboost.png}
  \caption{F1-score per lettera del best model.\newline
  \textit{Cosa mostra:} performance della migliore configurazione su ciascuna lettera. \textit{Cosa ci dice:} individua classi forti e classi deboli su cui concentrare miglioramenti mirati.}
\end{figure}

\section{Discussione}
I risultati mostrano in modo coerente che:
\begin{itemize}[leftmargin=1.5em]
  \item il feature engineering morfologico produce prestazioni competitive senza ricorrere a CNN;
  \item il setup subject-wise riduce in modo sostanziale il rischio di overfitting a uno specifico stile;
  \item le confusioni residue sono semanticamente plausibili (coppie graficamente simili).
\end{itemize}

Dal punto di vista metodologico, il valore principale è la tracciabilità della decisione:
si può collegare un errore a gruppi specifici di feature e quindi proporre miglioramenti mirati.
L'approccio è quindi particolarmente adatto quando si richiedono \textbf{interpretabilità}, \textbf{diagnostica} e \textbf{controllo} del processo.

\section{Conclusioni e sviluppi futuri}
Il progetto fornisce una pipeline completa, ripetibile e interpretabile per l'analisi e la classificazione di lettere manoscritte.
Il contributo non è solo nel valore numerico finale, ma nella costruzione di un processo leggibile in tutte le sue fasi: dalla morfologia all'analisi dell'errore.

Come sviluppi futuri:
\begin{enumerate}[leftmargin=1.5em]
  \item feature selection mirata sulle classi più deboli (riduzione ridondanza e maggiore robustezza);
  \item tuning iperparametri con protocollo nested CV per una stima ancora più rigorosa;
  \item estensione ibrida con backbone CNN leggero, mantenendo le feature morfologiche come canale interpretabile;
  \item analisi dell'importanza delle feature per gruppo, per guidare nuove iterazioni di design.
\end{enumerate}

\appendix
\section{Riproducibilità}
\subsection{Script principali}
\begin{itemize}[leftmargin=1.5em]
  \item \texttt{extract\_characters\_cvl.py}: segmenta caratteri da immagini CVL (single/batch), applica filtri geometrici/outlier e salva un manifest con bounding box.
  \item \texttt{analyze\_letter.py}: esegue preprocessing+morfologia e costruisce il dataset di feature campione per campione.
  \item \texttt{train\_noncnn\_advanced.py}: allena e confronta SVM/ExtraTrees/XGBoost su split holdout per soggetto.
  \item \texttt{crossval\_subjectwise.py}: stima robustezza con CV subject-wise a 5 fold.
  \item \texttt{ablation\_noncnn\_morphology.py}: studia l'impatto dei gruppi di feature (ablation).
\end{itemize}

\subsection{Output principali}
\begin{itemize}[leftmargin=1.5em]
  \item \texttt{output/extracted\_characters/extraction\_manifest.csv}
  \item \texttt{output/classification\_advanced/organized/01\_holdout/advanced\_summary.csv}
  \item \texttt{output/classification\_advanced/organized/01\_holdout/advanced\_metrics.json}
  \item \texttt{output/classification\_advanced/organized/02\_cv\_subjectwise/cv\_subjectwise\_summary.csv}
  \item \texttt{output/classification\_advanced/organized/03\_ablation/ablation\_summary.csv}
  \item \texttt{output/classification\_advanced/organized/03\_ablation/ablation\_group\_importance\_ranked.csv}
  \item \texttt{output/classification\_advanced/organized/03\_ablation/ablation\_feature\_group\_coverage.csv}
  \item \texttt{output/classification\_advanced/organized/03\_ablation/ablation\_delta\_macro\_f1.png}
  \item \texttt{output/classification\_advanced/organized/03\_ablation/ablation\_report.json}
  \item \texttt{output/classification\_advanced/organized/04\_diagnostics/confusion\_matrices/raw/confusion\_matrix\_xgboost.png}
  \item \texttt{output/classification\_advanced/organized/04\_diagnostics/confusion\_matrices/raw/confusion\_matrix\_svm\_rbf.png}
  \item \texttt{output/classification\_advanced/organized/04\_diagnostics/confusion\_matrices/raw/confusion\_matrix\_extra\_trees.png}
  \item \texttt{output/classification\_advanced/organized/04\_diagnostics/confusion\_matrices/normalized/confusion\_matrix\_xgboost\_normalized.png}
  \item \texttt{output/classification\_advanced/organized/04\_diagnostics/confusion\_matrices/normalized/confusion\_matrix\_svm\_rbf\_normalized.png}
  \item \texttt{output/classification\_advanced/organized/04\_diagnostics/confusion\_matrices/normalized/confusion\_matrix\_extra\_trees\_normalized.png}
  \item \texttt{output/classification\_advanced/organized/04\_diagnostics/feature\_importance/xgboost\_feature\_importance\_top20.png}
\end{itemize}

\subsection{Procedura minima di esecuzione}
\begin{enumerate}[leftmargin=1.5em]
  \item (Opzionale) Segmentazione caratteri CVL: \texttt{python extract\_characters\_cvl.py --input <path> --output <path> --save-debug --min-area 30 --min-width 5 --min-height 8 --max-aspect-ratio 10 --opening-kernel 2}
  \item Generare le feature: \texttt{python analyze\_letter.py}
  \item Allenare e valutare in holdout: \texttt{python train\_noncnn\_advanced.py}
  \item Eseguire CV subject-wise: \texttt{python crossval\_subjectwise.py}
  \item (Opzionale) Eseguire ablation: \texttt{python ablation\_noncnn\_morphology.py}
\end{enumerate}

\end{document}
